\documentclass[11pt]{aghdpl}
% \documentclass[en,11pt]{aghdpl}  % praca w języku angielskim

% Lista wszystkich języków stanowiących języki pozycji bibliograficznych użytych w pracy.
\usepackage[polish,english]{babel}
\usepackage[utf8]{inputenc}
\usepackage{mathtools}
\let\openbox\relax
\usepackage{amsthm}
\usepackage{amsfonts}


\usepackage{csquotes}
\DeclareQuoteAlias{croatian}{polish}

% --- < bibliografia > ---
\usepackage[
    style=numeric,
    sorting=none,
    language=autobib,
    autolang=other,
    urldate=iso,
    backref=false,
    isbn=true,
    url=false,
    maxbibnames=3,
    backend=biber
]{biblatex}

\addbibresource{bibliografia.bib}

\usepackage{caption}
\captionsetup[figure]{name=Rys.}
\captionsetup[table]{name=Tabela}

\usepackage{array}
\usepackage{tabularx}
\usepackage{multirow}
\usepackage{booktabs}
\usepackage{makecell}
\usepackage[flushleft]{threeparttable}

% defines the X column to use m (\parbox[c]) instead of p (`parbox[t]`)
\newcolumntype{C}[1]{>{\hsize=#1\hsize\centering\arraybackslash}X}


% --- < listingi > ---

\usepackage{courier}
\usepackage{listings}
\lstloadlanguages{TeX}

\lstset{
	literate={ą}{{\k{a}}}1
           {ć}{{\'c}}1
           {ę}{{\k{e}}}1
           {ó}{{\'o}}1
           {ń}{{\'n}}1
           {ł}{{\l{}}}1
           {ś}{{\'s}}1
           {ź}{{\'z}}1
           {ż}{{\.z}}1
           {Ą}{{\k{A}}}1
           {Ć}{{\'C}}1
           {Ę}{{\k{E}}}1
           {Ó}{{\'O}}1
           {Ń}{{\'N}}1
           {Ł}{{\L{}}}1
           {Ś}{{\'S}}1
           {Ź}{{\'Z}}1
           {Ż}{{\.Z}}1,
	basicstyle=\footnotesize\ttfamily,
}


%---------------------------------------------------------------------------

\author{Marcin Szpyrka}
\shortauthor{M. Szpyrka}

%\titlePL{Przygotowanie bardzo długiej i pasjonującej pracy dyplomowej w~systemie~\LaTeX}
%\titleEN{Preparation of a very long and fascinating bachelor or master thesis in \LaTeX}

\titlePL{Przygotowanie pracy dyplomowej w~systemie~\LaTeX}
\titleEN{Thesis in \LaTeX} % Podajemy tytuł w obu językach


\shorttitlePL{Przygotowanie pracy dyplomowej w~systemie \LaTeX} % skrócona wersja tytułu jeśli jest bardzo długi, MUSI BYĆ ZDEFINIOWANY
\shorttitleEN{Preparation of a long and fascinating thesis in \LaTeX}

\thesistype{Praca dyplomowa}
%\thesistype{Master of Science Thesis}

\supervisor{prof. dr hab. Marcin Szpyrka}
%\supervisor{Marcin Szpyrka PhD, DSc}

\degreeprogramme{Informatyka}
%\degreeprogramme{Computer Science}

\date{2022}

\department{Katedra Informatyki Stosowanej}
%\department{Department of Applied Computer Science}

\faculty{Wydział Elektrotechniki, Automatyki,\protect\\[-1mm] Informatyki i Inżynierii Biomedycznej}
%\faculty{Faculty of Electrical Engineering, Automatics, Computer Science and Biomedical Engineering}

\acknowledgements{Serdecznie dziękuję \dots tu ciąg dalszych podziękowań np. dla promotora, żony, sąsiada itp.}

\makeatletter
\def\thesisheaders{
\fancyfoot[LE,RO]{\small \@shortauthor\quad\textit{\@shorttitleEN}}
}
\def\thesistable{
\begin{tabular}{p{45mm}l}
Author: & {\itshape \@author}\\[-1mm]
Degree programme: & {\itshape \@degreeprogramme}\\[-1mm]
Supervisor: & {\itshape \@supervisor}\\
\end{tabular}
}
\makeatother
\thesisheaders

\setlength{\cftsecnumwidth}{10mm}

%---------------------------------------------------------------------------
\setcounter{secnumdepth}{4}
\brokenpenalty=10000\relax

% rubber: bibtex.path ./

\begin{document}

\titlepages

% Ponowne zdefiniowanie stylu `plain`, aby usunąć numer strony z pierwszej strony spisu treści i poszczególnych rozdziałów.
\fancypagestyle{plain}
{
	% Usuń nagłówek i stopkę
	\fancyhf{}
	% Usuń linie.
	\renewcommand{\headrulewidth}{0pt}
	\renewcommand{\footrulewidth}{0pt}
}

\setcounter{tocdepth}{2}
\tableofcontents
\clearpage

\chapter{Introduction}
\label{cha:introduction}

\section{Background}
% TODO

\section{Problem Statement}
% TODO

\section{Research Motivation}
% TODO

\section{Thesis Objectives}
% TODO

\section{Scope of the Thesis}
% TODO

\section{Thesis Structure}
% TODO

\chapter{Background and Related Work}
\label{cha:background-and-related-work}

\section{Scholarly Search Systems}
% TODO

\section{AI-Assisted Literature Retrieval}
% TODO

\section{Bias in Information Retrieval}
% TODO

\section{Metadata Enrichment in Scholarly Databases}
% TODO

\section{Existing Evaluation Approaches}
% TODO

\section{Research Gap}
% TODO

\chapter{Research Objectives and Methodology}
\label{cha:research-objectives-and-methodology}

\section{Research Goals}
% TODO

\section{Research Questions}
% TODO

\section{Scope and Assumptions}
% TODO

\section{Comparative Evaluation Methodology}
% TODO

\section{Bias-Related Metrics}
% TODO

\section{Limitations of the Methodology}
% TODO

\chapter{System Design and Architecture}
\label{cha:system-design-and-architecture}

\section{System Overview}
% TODO

\section{Backend Architecture}
% TODO

\section{Frontend Architecture}
% TODO

\section{Data Flow and Processing Pipeline}
% TODO

\section{Artifact Generation and Reproducibility}
% TODO

\section{Design Decisions}
% TODO

\chapter{Implementation}
\label{cha:implementation}

\section{Connector Layer}
% TODO

\section{Enrichment Module}
% TODO

\section{Metrics Computation}
% TODO

\section{Reporting and Visualization}
% TODO

\section{Run Management and Persistence}
% TODO

\section{User Interface Implementation}
% TODO

\chapter{Evaluation and Results}
\label{cha:evaluation-and-results}

\section{Experimental Setup}
% TODO

\section{Query Selection}
% TODO

\section{Tested Platforms and Models}
% TODO

\section{Metadata Coverage and Completeness}
% TODO

\section{Comparative Results}
% TODO

\section{Bias-Related Observations}
% TODO

\section{Summary of Findings}
% TODO

\chapter{Discussion}
\label{cha:discussion}

\section{Interpretation of Findings}
% TODO

\section{Strengths of the Proposed System}
% TODO

\section{Limitations and Threats to Validity}
% TODO

\section{Practical Usefulness}
% TODO

\section{Lessons Learned}
% TODO

\chapter{Conclusions and Future Work}
\label{cha:conclusions-and-future-work}

\section{Thesis Summary}
% TODO

\section{Main Contributions}
% TODO

\section{Answers to the Research Questions}
% TODO

\section{Future Work}
% TODO




% itd.
% \appendix
% \include{dodatekA}
% \include{dodatekB}
% itd.

\printbibliography

\end{document}
